\documentclass[11pt,letterpaper]{article}

\usepackage[spanish,es-noshorthands]{babel}
\usepackage{fontspec}
\setmainfont{Source Serif Pro}
\usepackage{microtype}
\usepackage{geometry}
\usepackage{booktabs}
\usepackage{tabularx}
\usepackage{multirow}
\usepackage{enumitem}
\usepackage{xcolor}
\usepackage{graphicx}
\usepackage{csquotes}
\usepackage{amsmath}
\usepackage{siunitx}
\usepackage[version=4]{mhchem}
\usepackage{tikz}
\usetikzlibrary{arrows.meta,positioning,shapes.geometric,fit,backgrounds,
                patterns,decorations.pathreplacing,calc}
\usepackage{placeins}
\usepackage{xurl}
\usepackage[hidelinks]{hyperref}
\usepackage[backend=biber,style=apa,sorting=nyt]{biblatex}
\DeclareLanguageMapping{spanish}{spanish-apa}

\geometry{margin=2.3cm}
% Los rótulos de las figuras usan nodos de ancho fijo con texto alineado, y la
% bibliografía compone en modo sloppy: ambos generan avisos de línea holgada
% que son cosméticos. Se silencian para que `make check` siga siendo sensible
% a los desbordes reales.
\hbadness=10000
\setlength{\emergencystretch}{1em}
\setlength{\parindent}{0pt}
\setlength{\parskip}{0.55em}
\setlist{nosep}
\renewcommand{\topfraction}{0.92}
\renewcommand{\bottomfraction}{0.85}
\renewcommand{\textfraction}{0.06}
\renewcommand{\floatpagefraction}{0.72}
\setcounter{topnumber}{2}
\setcounter{bottomnumber}{2}
\setcounter{totalnumber}{3}
\addbibresource{referencias.bib}
\AtBeginBibliography{\sloppy}

\sisetup{
  output-decimal-marker = {,},
  group-separator = {.},
  group-digits = integer,
  group-minimum-digits = 4,
  detect-all
}
\DeclareSIUnit{\pkm}{pax\text{-}km}
\DeclareSIUnit{\pax}{pax}
\DeclareSIUnit{\litro}{L}

% Colores: dato de partida, resultado calculado, magnitud despreciada,
% alerta o incoherencia, referencia externa y contexto neutro
\definecolor{cDato}{RGB}{70,90,112}
\definecolor{cCalc}{RGB}{16,110,104}
\definecolor{cAviso}{RGB}{176,42,42}
\definecolor{cNeutro}{RGB}{132,124,112}
\definecolor{cVerde}{RGB}{74,116,60}
\definecolor{cNaranja}{RGB}{186,98,52}

% --- Figura 2: fracción de combustible frente al alcance --------------------
\newcommand{\cRx}[1]{#1*0.655}     % miles de km -> cm (20 = 13,1 cm)
\newcommand{\cRy}[1]{#1*0.0620}     % %   -> cm  (100 %    =  6,2 cm)

% --- Figura 4: cascada del combustible --------------------------------------
\newcommand{\cWy}[1]{#1*0.00178}    % kg  -> cm  (3.000 kg =  5,3 cm)
% #1 x izquierda, #2 base, #3 cima, #4 color, #5 rótulo, #6 valor,
% #7 altura del rótulo (los rótulos se alternan en dos filas)
\newcommand{\tramo}[7]{%
  \fill[#4] (#1,{\cWy{#2}}) rectangle ({#1+1.14},{\cWy{#3}});%
  \node[above, font=\scriptsize\bfseries, text=#4, fill=white,
        inner sep=1.3pt] at ({#1+0.57},{\cWy{#3}}) {#6};%
  \draw[cNeutro!45] ({#1+0.57},-0.06) -- ({#1+0.57},{#7+0.10});%
  \node[below, font=\scriptsize, text width=2.35cm, align=center,
        text=cNeutro!80!black] at ({#1+0.57},#7) {#5};%
}

% --- Figura 5: tornado de sensibilidad --------------------------------------
\newcommand{\tSx}[1]{(#1-1700)*0.0094}   % kg -> cm (1.700 a 3.000 = 12,2 cm)
% #1 y, #2 kg extremo bajo, #3 kg extremo alto, #4 rótulo, #5 texto bajo,
% #6 texto alto
\newcommand{\tornado}[6]{%
  \node[left, font=\scriptsize, text width=3.15cm, align=right]
       at (-0.52,#1) {#4};%
  \fill[cCalc!70] ({\tSx{#2}},#1-0.27) rectangle ({\tSx{2248}},#1+0.27);%
  \fill[cAviso!70] ({\tSx{2248}},#1-0.27) rectangle ({\tSx{#3}},#1+0.27);%
  \node[left, font=\scriptsize, text=cCalc!80!black]
       at ({\tSx{#2}-0.10},#1) {#5};%
  \node[right, font=\scriptsize, text=cAviso]
       at ({\tSx{#3}+0.10},#1) {#6};%
}

\title{\vspace{-1.3cm}\textbf{¿Cuánto queroseno hace falta para volar de
Boston a Los Ángeles?}\\
\large Estimación mediante la ecuación de alcance y resistencia de Bréguet,
comprobación de la estimación y análisis de por qué el cálculo es difícil}
\author{Carlos Luis Rojas Aragonés\\
\small Gestión del Desarrollo Tecnológico}
\date{\small 7 de septiembre de 2026}

\begin{document}

\maketitle
\vspace{-0.8em}

\section{El resultado, antes de justificarlo}

Con los datos del enunciado, la ecuación de alcance de Bréguet da
\textbf{\SI{2248}{\kilogram} de queroseno}, es decir unas \textbf{2,25
toneladas} o \textbf{\SI{2800}{\litro}}, para recorrer los
\SI{5000}{\kilo\meter} en crucero. Eso es el \textbf{22,5 por ciento} del peso
al despegue del avión: casi una cuarta parte de lo que sale del suelo es
combustible que se va a quemar por el camino.

La cifra sale de una cadena de conversión muy corta. El queroseno lleva
\SI{42.8}{\mega\joule} de energía química por kilogramo; el conjunto
motor más hélice o tobera entrega al aire el 30 por ciento de eso en forma de
trabajo propulsivo; y el ala convierte cada newton de arrastre vencido en
quince newtons de peso sostenido. El producto de esas tres cosas es una
distancia, y esa distancia es la respuesta.

El Cuadro~\ref{cuadro:respuesta} adelanta el resultado completo, incluidas las
dos lecturas posibles del enunciado y lo que le falta a la cifra para ser un
plan de vuelo.

\begin{table}[htbp]
\small
\caption{La respuesta, en las cuatro formas en que puede pedirse, y lo que hay
que añadirle para que sea utilizable. Cálculo propio con los datos del
enunciado y \(h=\SI{42.8}{\mega\joule\per\kilogram}\).}
\label{cuadro:respuesta}
\begin{tabularx}{\textwidth}{@{}Xrl@{}}
\toprule
\textbf{Magnitud} & \textbf{Valor} & \textbf{Procedencia} \\
\midrule
Parámetro adimensional \(x = Rg/[\eta_o h (L/D)]\) & \num{0.2547} &
  ecuación~\eqref{eq:x} \\
Razón de pesos \(W_i/W_f = e^{x}\) & \num{1.2900} & ecuación~\eqref{eq:breguet} \\
\midrule
\textbf{Queroseno, lectura A} (las \SI{10}{\tonne} son el peso al despegue)
  & \textbf{\SI{2248}{\kilogram}} & \(m_i(1-e^{-x})\) \\
Queroseno, lectura B (las \SI{10}{\tonne} son el peso sin combustible)
  & \SI{2900}{\kilogram} & \(m_f(e^{x}-1)\) \\
Queroseno con la forma lineal de la ecuación (lectura A)
  & \SI{2547}{\kilogram} & ecuación~\eqref{eq:lineal} \\
\midrule
Volumen a \SI{0.804}{\kilogram\per\litro} (lectura A) & \SI{2796}{\litro} &
  \(m_f/\rho_{\text{comb}}\) \\
Energía química embarcada & \SI{96.2}{\giga\joule} & \(m_f h\) \\
\ce{CO2} emitido, factor \SI{3.16}{\kilogram\per\kilogram} & \SI{7.10}{\tonne} &
  \textcite{icaocalc} \\
Duración del crucero (resistencia) & \SI{4.63}{\hour} & \(R/V\) \\
\midrule
Combustible que habría que embarcar de verdad & \(\approx\)
  \SI{2850}{\kilogram} & Figura~\ref{fig:cascada} \\
Combustible si se corrige el \(L/D\) a la condición real de vuelo &
  \(>\SI{4900}{\kilogram}\) & Cuadro~\ref{cuadro:coherencia} \\
\bottomrule
\end{tabularx}
\end{table}

Las dos últimas filas son el motivo de que la pregunta final del enunciado
(\enquote{¿por qué es un cálculo complicado?}) tenga una respuesta larga y no
una respuesta retórica. La ecuación se resuelve en cinco líneas. El problema
que la ecuación pretende resolver no se resuelve en cinco líneas, y la
distancia entre las dos cosas es medible: la doy cuantificada en la
sección~\ref{sec:complicado}.

\section{La ecuación: de dónde sale y qué dice}

\subsection{Nota sobre el nombre}

Conviene aclarar dos cosas del enunciado antes de usarlo. La primera es que
\enquote{resistencia}, en \enquote{ecuación de alcance y resistencia},
traduce el inglés \emph{endurance} y significa \textbf{tiempo que el avión
puede mantenerse en el aire}, no \emph{drag}, que en español es arrastre. Son
dos magnitudes distintas y las dos aparecen en este problema: el arrastre es
la fuerza que hay que vencer y la resistencia es la duración del vuelo,
\SI{4.63}{\hour} en este caso. La segunda es que la ecuación se conoce por el
nombre de Louis Charles Bréguet, que la publicó en 1923, pero la habían
derivado antes Devillers en 1918 y \textcite{coffin1920} en el informe 69 de
la NACA. El nombre es una convención, no una atribución.

\subsection{Derivación completa}

La ecuación no es una correlación empírica: se deduce de tres afirmaciones
elementales. Sigo la derivación de \textcite{mitunified}, que es la que usa el
módulo, y la de \textcite{torenbeek2009}.

\textbf{Primera afirmación: el vuelo es recto, horizontal y uniforme.} Las
cuatro fuerzas se equilibran dos a dos, como en el Panel A de la
Figura~\ref{fig:cadena}:
\begin{equation}
L = W, \qquad T = D = \frac{L}{L/D} = \frac{W}{L/D}.
\label{eq:equilibrio}
\end{equation}
Esta es la única aparición de la aerodinámica en todo el desarrollo, y entra
por un solo número, \(L/D\). El enunciado nos lo da hecho: 15.

\textbf{Segunda afirmación: el avión adelgaza a medida que quema.} La tasa de
cambio del peso es el flujo másico de combustible multiplicado por la gravedad,
con signo negativo:
\begin{equation}
\frac{dW}{dt} = -\dot{m}_f\, g .
\label{eq:masa}
\end{equation}

\textbf{Tercera afirmación: definición de eficiencia total.} Lo que se obtiene
es potencia propulsiva; lo que se paga es potencia química:
\begin{equation}
\eta_o \equiv \frac{\text{lo que se obtiene}}{\text{lo que se paga}}
        = \frac{\text{potencia propulsiva}}{\text{potencia de combustible}}
        = \frac{T\,V}{\dot{m}_f\, h}
\quad\Longrightarrow\quad
\dot{m}_f = \frac{T\,V}{\eta_o\, h}.
\label{eq:eta}
\end{equation}
El enunciado nos da \(\eta_o = 0{,}3\). Ese número es el producto de un
rendimiento termodinámico (lo que el ciclo del motor convierte en energía
cinética del chorro) por un rendimiento propulsivo (lo que ese chorro convierte
en empuje útil), y 0,3 es un valor razonable para un turbofán de generación
media \parencite{mitthermo}.

\textbf{Combinación e integración.} Sustituyendo \eqref{eq:equilibrio} y
\eqref{eq:eta} en \eqref{eq:masa} se obtiene una ecuación diferencial en la que
ya no aparece ni el empuje ni el flujo de combustible:
\begin{equation}
\frac{dW}{dt}
  = -\frac{T V g}{\eta_o h}
  = -\,\frac{W}{L/D}\cdot\frac{V g}{\eta_o h}
\quad\Longrightarrow\quad
\frac{dW}{W} = -\,\frac{V g}{\eta_o\, h\,(L/D)}\; dt .
\label{eq:edo}
\end{equation}
El paso decisivo es el de la izquierda: \textbf{la variable dependiente aparece
en el numerador y en el denominador}. Eso es lo que convierte el problema en
logarítmico y no en proporcional, y es la raíz de casi todas las dificultades
que discuto en la sección~\ref{sec:complicado}. Integrando entre el estado
inicial y el final, con \(\eta_o\), \(L/D\) y \(V\) constantes:
\begin{equation}
\ln\frac{W_i}{W_f} = \frac{V g\, t}{\eta_o\, h\,(L/D)} .
\end{equation}

De ahí salen las dos formas que pide el enunciado. Multiplicando por la
velocidad, con \(R = V t\), la velocidad se cancela y queda el
\textbf{alcance}:
\begin{equation}
\boxed{\;R = \eta_o\,\frac{h}{g}\,\frac{L}{D}\,\ln\frac{W_i}{W_f}\;}
\label{eq:breguet}
\end{equation}
y sin multiplicar por la velocidad queda la \textbf{resistencia}, es decir la
duración:
\begin{equation}
E = \frac{\eta_o}{V}\,\frac{h}{g}\,\frac{L}{D}\,\ln\frac{W_i}{W_f}
  = \frac{R}{V} .
\label{eq:endurance}
\end{equation}

\begin{figure}[htbp]
\centering
\begin{tikzpicture}

% ================= Panel A: equilibrio de fuerzas =================
\node[anchor=south west, font=\small\bfseries, text width=6.6cm, align=left]
     at (-0.4,3.30)
     {Panel A. Las cuatro fuerzas en crucero, con los valores de este
      problema};

% silueta de perfil
\begin{scope}[shift={(3.2,0.60)}]
  \fill[cNeutro!35, rounded corners=6pt] (-1.30,-0.17) rectangle (1.05,0.25);
  \fill[cNeutro!35] (1.00,-0.10) -- (1.62,0.04) -- (1.00,0.22) -- cycle;
  \fill[cDato!45] (-1.28,0.22) -- (-0.86,0.22) -- (-1.02,0.92) --
                  (-1.24,0.92) -- cycle;
  \fill[cDato!45] (-1.30,0.20) -- (-0.72,0.26) -- (-0.72,0.32) --
                  (-1.30,0.30) -- cycle;
  \fill[cDato!60] (-0.30,-0.06) -- (0.52,0.02) -- (0.52,0.09) --
                  (-0.30,0.04) -- cycle;
  \fill[cNeutro!55, rounded corners=3pt] (-0.16,-0.42) rectangle (0.46,-0.14);
\end{scope}

% sustentación
\draw[cCalc, very thick, -{Stealth[length=6pt]}] (3.20,1.05) -- (3.20,2.95);
\node[right, font=\scriptsize, text=cCalc, align=left] at (3.32,2.66)
     {$L = W = \SI{98100}{\newton}$};
% peso
\draw[cAviso, very thick, -{Stealth[length=6pt]}] (3.20,0.12) -- (3.20,-1.55);
\node[right, font=\scriptsize, text=cAviso, align=left] at (3.32,-1.28)
     {$W = m\,g$, \emph{decreciente}};
% empuje
\draw[cNaranja, very thick, -{Stealth[length=6pt]}] (5.00,0.70) -- (6.30,0.70);
\node[above, font=\scriptsize, text=cNaranja] at (5.65,0.78) {$T$};
% arrastre
\draw[cDato, very thick, -{Stealth[length=6pt]}] (1.72,0.70) -- (0.42,0.70);
\node[above, font=\scriptsize, text=cDato] at (1.07,0.80) {$D$};
\node[below, font=\scriptsize, text=cDato] at (1.07,0.58) {$\dfrac{W}{L/D}$};
\node[font=\scriptsize, align=center, text=cNeutro!85!black] at (1.07,-0.62)
     {$\SI{6540}{\newton}$ al inicio\\$\SI{5070}{\newton}$ al final};

% ================= Panel B: las cuatro disciplinas =================
\begin{scope}[shift={(8.75,0.55)}]
\node[anchor=south west, font=\small\bfseries, text width=6.6cm, align=left]
     at (-0.55,2.75)
     {Panel B. Los cuatro factores y la disciplina de la que sale cada uno};

\node[font=\normalsize] (r)  at (0,1.45) {$R \;=\;$};
\node[font=\normalsize, text=cNaranja, right=1pt of r]  (f1) {$\eta_o$};
\node[font=\normalsize, right=4pt of f1] (x1) {$\times$};
\node[font=\normalsize, text=cVerde, right=4pt of x1]   (f2) {$\dfrac{h}{g}$};
\node[font=\normalsize, right=4pt of f2] (x2) {$\times$};
\node[font=\normalsize, text=cDato, right=4pt of x2]    (f3) {$\dfrac{L}{D}$};
\node[font=\normalsize, right=4pt of f3] (x3) {$\times$};
\node[font=\normalsize, text=cAviso, right=4pt of x3]
     (f4) {$\ln\dfrac{W_i}{W_f}$};

% los rótulos se alternan en dos alturas para que no se solapen
\draw[cNaranja, decorate, decoration={brace, mirror, amplitude=3.2pt}]
  (f1.south west) -- (f1.south east);
\node[below, font=\scriptsize, text=cNaranja, text width=1.9cm, align=center]
  at ($(f1.south)+(0,-0.14)$) {propulsión\\\num{0.30}};

\draw[cVerde, decorate, decoration={brace, mirror, amplitude=3.2pt}]
  (f2.south west) -- (f2.south east);
\draw[cVerde, thin] ($(f2.south)+(0,-0.12)$) -- ($(f2.south)+(0,-0.94)$);
\node[below, font=\scriptsize, text=cVerde, text width=2.1cm, align=center]
  at ($(f2.south)+(0,-0.94)$) {química del combustible\\\SI{4363}{\kilo\meter}};

\draw[cDato, decorate, decoration={brace, mirror, amplitude=3.2pt}]
  (f3.south west) -- (f3.south east);
\node[below, font=\scriptsize, text=cDato, text width=1.9cm, align=center]
  at ($(f3.south)+(0,-0.14)$) {aerodinámica\\\num{15}};

\draw[cAviso, decorate, decoration={brace, mirror, amplitude=3.2pt}]
  (f4.south west) -- (f4.south east);
\draw[cAviso, thin] ($(f4.south)+(0,-0.12)$) -- ($(f4.south)+(0,-0.94)$);
\node[below, font=\scriptsize, text=cAviso, text width=2.3cm, align=center]
  at ($(f4.south)+(0,-0.94)$) {estructuras y carga de pago\\\num{0.2547}};

\node[font=\scriptsize, align=left, text width=6.6cm, text=cNeutro!85!black]
     at (2.7,-2.10)
     {El producto de los tres primeros factores es una distancia,
      \SI{19633}{\kilo\meter}. El cuarto dice qué fracción de esa distancia se
      recorre de verdad.};
\end{scope}

\end{tikzpicture}
\caption{La ecuación de Bréguet en dos vistas. Panel A: el equilibrio de
fuerzas del que se deduce, con los valores numéricos de este problema. El peso
es la única de las cuatro fuerzas que cambia durante el vuelo, y de ahí sale el
logaritmo. Panel B: los cuatro factores del segundo miembro de la
ecuación~\eqref{eq:breguet}, cada uno perteneciente a una disciplina distinta.
Esta descomposición es la razón por la que \textcite{mitunified} usa esta
ecuación como puerta de entrada a la ingeniería aeronáutica: es el punto en el
que propulsión, aerodinámica y estructuras se multiplican. Elaboración propia
a partir de \textcite{mitunified,deweck2022}.}
\label{fig:cadena}
\end{figure}

\subsection{La forma lineal y por qué la menciono}

El módulo, siguiendo a \textcite{tennekes2009}, presenta también una versión
simplificada que se obtiene suponiendo que todo el combustible se consume en la
misma condición de vuelo y que el peso apenas cambia:
\begin{equation}
R \;\approx\; \eta_o\,\frac{h}{g}\,\frac{L}{D}\cdot\frac{W_f}{W},
\label{eq:lineal}
\end{equation}
donde \(W_f\) es el peso del combustible. No es otra ecuación: es el primer
término del desarrollo de la ecuación~\eqref{eq:breguet}, porque
\(\ln(1/(1-\zeta)) \approx \zeta\) cuando la fracción de combustible \(\zeta\)
es pequeña. La menciono por dos motivos. El primero es que da una respuesta
distinta, \SI{2547}{\kilogram} en lugar de \SI{2248}{\kilogram}, un 13 por
ciento más, y conviene saber de dónde viene esa discrepancia antes de que
aparezca. El segundo es que la propia aproximación deja ver el mecanismo:
\textbf{cada kilogramo de combustible que se embarca hay que llevarlo también
en volandas hasta que se quema}, y la forma lineal es precisamente la que
ignora eso.

De la ecuación~\eqref{eq:lineal} sale además una observación que
\textcite{tennekes2009} subraya y que aquí resulta muy útil: el cociente
\(h/g\) tiene dimensiones de longitud. Para el queroseno,
\(h/g = \SI{42.8e6}{\joule\per\kilogram} / \SI{9.81}{\meter\per\second\squared}
= \SI{4363}{\kilo\meter}\). Ese es el alcance de referencia de un combustible,
y el alcance real de cualquier vehículo es del mismo orden de magnitud
multiplicado por \(\eta_o (L/D)\), que para este avión vale \num{4.5}. Antes de
calcular nada ya sabemos que \SI{5000}{\kilo\meter} es una fracción moderada de
lo alcanzable, y por tanto que la respuesta no va a ser ni el 2 ni el 80 por
ciento del peso del avión.

\section{Datos de partida}

El Cuadro~\ref{cuadro:datos} recoge todo lo que hace falta, en unidades del
sistema internacional, y declara explícitamente qué datos entran en la ecuación
y cuáles no.

\begin{table}[htbp]
\small
\caption{Datos de partida, convertidos a unidades coherentes. La última
columna indica si el dato entra en la ecuación~\eqref{eq:breguet} o si se usa
para otra cosa, que es una distinción que el enunciado no hace y que conviene
hacer.}
\label{cuadro:datos}
\begin{tabularx}{\textwidth}{@{}p{4.35cm}p{1.5cm}rXp{2.15cm}@{}}
\toprule
\textbf{Magnitud} & \textbf{Símbolo} & \textbf{Valor (SI)} &
\textbf{Origen} & \textbf{¿Entra en la ecuación?} \\
\midrule
Distancia Boston--Los Ángeles & \(R\) & \SI{5.0e6}{\meter} & enunciado &
  sí \\
Masa total al despegue & \(m_i\) & \SI{1.0e4}{\kilogram} & enunciado &
  sí \\
Relación sustentación/arrastre & \(L/D\) & \num{15} & enunciado & sí \\
Eficiencia total & \(\eta_o\) & \num{0.30} & enunciado & sí \\
Densidad energética del queroseno & \(h\) &
  \SI{42.8e6}{\joule\per\kilogram} & sección 3 del módulo;
  \textcite{tennekes2009} & sí \\
Aceleración de la gravedad & \(g\) &
  \SI{9.81}{\meter\per\second\squared} & constante & sí \\
\midrule
Velocidad de crucero & \(V\) & \SI{300}{\meter\per\second} & enunciado &
  \textbf{no} en el alcance; sí en la resistencia \\
Densidad del aire a \SI{6000}{\meter} & \(\rho\) &
  \SI{0.6125}{\kilogram\per\meter\cubed} & enunciado
  (\(\rho_0/2\), con \(\rho_0=\SI{1.225}{\kilogram\per\meter\cubed}\)) &
  \textbf{no}: sirve para comprobar la coherencia \\
\midrule
Densidad del queroseno & \(\rho_{\text{comb}}\) &
  \SI{804}{\kilogram\per\meter\cubed} & \textcite{chevron2007} &
  sólo para pasar a litros \\
Factor de emisión del queroseno & &
  \SI{3.16}{\kilogram\per\kilogram} de \ce{CO2} & \textcite{icaocalc} &
  sólo para la huella \\
\bottomrule
\end{tabularx}
\end{table}

\subsection{La densidad energética del queroseno}

Es el dato que el enunciado remite a la sección 3 del módulo. El valor que allí
se da, tomado de la tabla de poderes caloríficos de \textcite{tennekes2009}, es
\(h \approx \SI{42.5}{\mega\joule\per\kilogram}\) para el grupo
\enquote{queroseno, gasolina y gasóleo}, y la literatura de referencia sitúa el
poder calorífico inferior del Jet A-1 entre \num{42.8} y
\SI{43.2}{\mega\joule\per\kilogram} \parencite{chevron2007}. Uso
\SI{42.8}{\mega\joule\per\kilogram}, que es el valor mínimo de especificación
del combustible de aviación y el que emplea \textcite{mitunified} en el mismo
ejercicio.

Merece la pena detenerse un segundo en este dato, porque encierra la primera de
las trampas del problema. Lo que hay que usar es el \textbf{poder calorífico
inferior}, no el superior. La diferencia entre los dos es el calor latente del
vapor de agua que sale por la tobera: unos \SI{3.6}{\mega\joule\per\kilogram},
un 8 por ciento. Ese vapor se va a la atmósfera sin condensar, de modo que su
calor latente no está disponible para producir empuje. Si uno toma por
descuido el poder calorífico superior, el resultado sale un 7 por ciento bajo,
como se ve en el Cuadro~\ref{cuadro:sensibilidad}. Es un error pequeño en este
problema y grande en otros: para el hidrógeno la diferencia entre los dos
poderes caloríficos es del 18 por ciento.

\section{El cálculo, paso a paso}

\subsection{Los cinco pasos}

\textbf{Paso 1. El alcance de referencia del combustible.}
\begin{equation}
\frac{h}{g} = \frac{\SI{42.8e6}{\joule\per\kilogram}}
                   {\SI{9.81}{\meter\per\second\squared}}
            = \SI{4.363e6}{\meter} = \SI{4363}{\kilo\meter}.
\end{equation}

\textbf{Paso 2. El alcance característico del avión}, es decir el que
recorrería si todo su peso fuera combustible y pudiera quemarlo entero:
\begin{equation}
R^{*} \equiv \eta_o\,\frac{h}{g}\,\frac{L}{D}
      = 0{,}30 \times \SI{4363}{\kilo\meter} \times 15
      = \SI{19633}{\kilo\meter}.
\label{eq:rest}
\end{equation}
Esta cifra es útil por sí misma: es media vuelta al mundo, lo que dice que un
avión con estos números está lejos de poder circunnavegar el planeta pero muy
por encima de lo que necesita para cruzar los Estados Unidos.

\textbf{Paso 3. El parámetro adimensional.} La ecuación~\eqref{eq:breguet} se
reescribe como \(R = R^{*}\ln(W_i/W_f)\), de modo que el exponente que gobierna
todo es el cociente entre la distancia pedida y el alcance característico:
\begin{equation}
x \equiv \frac{R}{R^{*}} = \frac{R\,g}{\eta_o\, h\,(L/D)}
  = \frac{\SI{5.0e6}{\meter}\times\SI{9.81}{\meter\per\second\squared}}
         {0{,}30\times\SI{42.8e6}{\joule\per\kilogram}\times 15}
  = \frac{\num{4.905e7}}{\num{1.926e8}} = \num{0.2547}.
\label{eq:x}
\end{equation}

\textbf{Paso 4. La razón de pesos.}
\begin{equation}
\frac{W_i}{W_f} = \frac{m_i}{m_f} = e^{x} = e^{0{,}2547} = \num{1.2900}.
\end{equation}

\textbf{Paso 5. La masa de combustible.} Aquí el enunciado admite dos lecturas,
y las doy las dos porque la diferencia entre ellas no es pequeña.

\emph{Lectura A: las \SI{10}{\tonne} son el peso total al despegue}, que es lo
que dice literalmente el enunciado (\enquote{este es el peso total}). Entonces
\(m_i = \SI{10000}{\kilogram}\) y el combustible es la diferencia entre el peso
inicial y el final:
\begin{align}
m_{\text{final}} &= \frac{m_i}{e^{x}}
  = \frac{\SI{10000}{\kilogram}}{1{,}2900} = \SI{7752}{\kilogram},\\
m_{\text{comb}}  &= m_i\left(1-e^{-x}\right)
  = \SI{10000}{\kilogram}\times(1-0{,}7752)
  = \boxed{\SI{2248}{\kilogram}} .
\label{eq:resultado}
\end{align}

\emph{Lectura B: las \SI{10}{\tonne} son el peso sin combustible}, es decir
estructura más pasajeros más carga, que es la interpretación que sugiere la
aclaración \enquote{incluye pasajeros y carga}. Entonces
\(m_{\text{final}} = \SI{10000}{\kilogram}\) y
\begin{equation}
m_{\text{comb}} = m_{\text{final}}\left(e^{x}-1\right)
  = \SI{10000}{\kilogram}\times 0{,}2900 = \SI{2900}{\kilogram},
\end{equation}
con un peso al despegue de \SI{12900}{\kilogram}. Las dos cifras difieren en un
29 por ciento, y la ambigüedad de una sola frase del enunciado es la mayor
fuente de incertidumbre de todo el ejercicio, por encima de la del valor de
\(h\) y comparable a la del \(L/D\). En lo que sigue uso la lectura A, que es
la literal, y arrastro la B en el análisis de sensibilidad.

\subsection{Qué forma tiene la respuesta}

La Figura~\ref{fig:curva} sitúa el resultado en la curva completa. Es la parte
del ejercicio que a mí más me ha servido, porque explica de un golpe tres cosas
que por separado parecen desconectadas: por qué la forma lineal se equivoca
poco aquí y mucho en un vuelo transpacífico, por qué existe un alcance máximo
aunque el depósito fuera infinito, y por qué el reactor de negocios de este
problema no llega a Los Ángeles con las reservas puestas.

\begin{figure}[htbp]
\centering
\begin{tikzpicture}

\node[anchor=south west, font=\small\bfseries] at (-0.9,6.65)
     {Fracción del peso al despegue que hay que llevar en queroseno, frente al
      alcance};

% rejilla y ejes
\foreach \v in {20,40,60,80,100} {
  \draw[cNeutro!28] (0,{\cRy{\v}}) -- (13.5,{\cRy{\v}});
  \node[left, font=\scriptsize, text=cNeutro] at (-0.06,{\cRy{\v}}) {\v\,\%};
}
\draw[cNeutro!70] (0,0) -- (13.5,0);
\draw[cNeutro!70] (0,0) -- (0,6.45);
\foreach \a in {0,2.5,5,7.5,10,12.5,15,17.5,20} {
  \draw[cNeutro!70] ({\cRx{\a}},0) -- ({\cRx{\a}},-0.09);
}
\foreach \a/\t in {0/0, 5/5.000, 10/10.000, 15/15.000, 20/20.000} {
  \node[below, font=\scriptsize, text=cNeutro] at ({\cRx{\a}},-0.11) {\t};
}
\node[below, font=\scriptsize, text=cNeutro] at (6.75,-0.62)
     {alcance en crucero (km)};

% asíntota del 100 %
\draw[cNeutro!75, dashed] (0,{\cRy{100}}) -- (13.5,{\cRy{100}});

% curva lineal (aproximación de Tennekes)
\draw[cNaranja, thick, dashed, samples=60, domain=0:19.633, variable=\r]
  plot ({\cRx{\r}},{\cRy{100*\r/19.633}});
\node[right, font=\scriptsize, text=cNaranja, align=left, rotate=27]
     at ({\cRx{12.1}},{\cRy{66}}) {forma lineal: $R/R^{*}$};

% curva exacta
\draw[cCalc, very thick, samples=90, domain=0:20, variable=\r]
  plot ({\cRx{\r}},{\cRy{100*(1-exp(-\r/19.633))}});
\node[right, font=\scriptsize, text=cCalc, align=left]
     at ({\cRx{15.4}},{\cRy{45}}) {forma exacta:\\$1-e^{-R/R^{*}}$};

% punto de operación
\draw[cDato, dashed] ({\cRx{5}},0) -- ({\cRx{5}},{\cRy{100}});
\fill[cCalc] ({\cRx{5}},{\cRy{22.48}}) circle (2.6pt);
\fill[cNaranja] ({\cRx{5}},{\cRy{25.47}}) circle (2.6pt);
\node[right, font=\scriptsize, text=cCalc, align=left]
     at ({\cRx{5.25}},{\cRy{17.5}})
     {\textbf{este problema}\\\SI{22.48}{\percent} = \SI{2248}{\kilogram}};
\node[left, font=\scriptsize, text=cNaranja, align=right]
     at ({\cRx{4.75}},{\cRy{32}})
     {\SI{25.47}{\percent} = \SI{2547}{\kilogram}\\(forma lineal,
      \num[retain-explicit-plus]{+13}\,\%)};

% R*
\draw[cAviso, dotted, thick] ({\cRx{19.633}},0) -- ({\cRx{19.633}},{\cRy{100}});
\node[above left, font=\scriptsize, text=cAviso, align=right]
     at ({\cRx{19.633}},{\cRy{101}})
     {$R^{*} = \SI{19633}{\kilo\meter}$};
\fill[cAviso] ({\cRx{19.633}},{\cRy{63.21}}) circle (2.2pt);
\node[left, font=\scriptsize, text=cAviso] at ({\cRx{19.4}},{\cRy{63.21}})
     {\SI{63.2}{\percent}};

% banda de capacidad de un reactor real de 10 t
\fill[cDato!16] (0,{\cRy{26}}) rectangle (13.5,{\cRy{28.5}});
\node[right, font=\scriptsize, text=cDato, align=left]
     at ({\cRx{8.05}},{\cRy{11.5}})
     {capacidad de depósito de un reactor real de \SI{10}{\tonne}:
      \SI{2750}{\kilogram} sobre \SI{9752}{\kilogram},
      el \SI{28.2}{\percent}};
\draw[cDato, -{Stealth[length=4pt]}] ({\cRx{10.4}},{\cRy{16.5}})
  -- ({\cRx{10.4}},{\cRy{25.2}});

% zona de validez de la forma lineal
\draw[cNeutro!85, |-|] ({\cRx{0.2}},{\cRy{92}}) -- ({\cRx{4}},{\cRy{92}});
\node[right, font=\scriptsize, text=cNeutro!85!black, align=left]
     at ({\cRx{4.15}},{\cRy{92}})
     {por debajo de $R \approx R^{*}/5$ las dos formas coinciden dentro
      del \SI{10}{\percent}};

\end{tikzpicture}
\caption{El resultado en su contexto. La curva verde es la
ecuación~\eqref{eq:breguet} resuelta para la fracción de combustible; la
naranja discontinua es la aproximación lineal~\eqref{eq:lineal}, que sobrestima
y que diverge sin límite mientras la exacta se satura. El alcance
característico \(R^{*} = \eta_o (h/g)(L/D) = \SI{19633}{\kilo\meter}\) es la
escala natural del problema: a \(R = R^{*}\) haría falta el
\SI{63.2}{\percent} del peso al despegue en combustible, y ningún avión
convencional pasa del 45 o 50 por ciento, de modo que \(R^{*}\) no es un
alcance alcanzable sino una constante de normalización. La banda azul es la
capacidad real de depósito de un reactor de negocios de esta masa
\parencite{learjet75}: el punto de operación cae justo por debajo, que es la
razón de que un avión así llegue a Los Ángeles desde Boston sólo sobre el
papel. Elaboración propia.}
\label{fig:curva}
\end{figure}

Tres lecturas de la figura:

\begin{itemize}
\item \textbf{La curva se satura y la aproximación lineal no.} A
  \SI{5000}{\kilo\meter} la diferencia entre las dos es del 13 por ciento; a
  \SI{15000}{\kilo\meter} la forma lineal daría el 76 por ciento del peso al
  despegue en combustible frente al 53 por ciento real, un error del 43 por
  ciento. La aproximación es utilizable exactamente en el régimen de este
  problema y no más allá.
\item \textbf{Existe un techo estructural, no energético.} Como la fracción de
  combustible nunca puede llegar al 100 por ciento, la ecuación de Bréguet dice
  que el alcance está limitado por lo que la estructura permite embarcar. Con
  un 45 por ciento de fracción de combustible, que es un valor de récord para
  un avión de transporte, este avión llegaría a \SI{11700}{\kilo\meter}.
\item \textbf{El punto de operación está incómodamente cerca del depósito.}
  Los \SI{2248}{\kilogram} son el 22,5 por ciento del peso al despegue, y el
  reactor de negocios real más parecido a este avión lleva
  \SI{2750}{\kilogram} de capacidad, el 28,2 por ciento
  \parencite{learjet75,aopalearjet}. Parece que sobra margen. La
  Figura~\ref{fig:cascada} muestra que no sobra.
\end{itemize}

\section{Cuatro comprobaciones independientes}

Un resultado de una línea conviene comprobarlo por caminos que no usen la
misma línea. Hago cuatro, resumidas en el Cuadro~\ref{cuadro:comprobaciones}.

\textbf{Comprobación 1: balance de energía con el peso medio.} El trabajo que
hay que hacer contra el arrastre es \(\overline{D}\,R\). El arrastre inicial es
\(D_i = W_i/(L/D) = \SI{98100}{\newton}/15 = \SI{6540}{\newton}\) y el final
\(D_f = \SI{5070}{\newton}\). Si se usa el peso inicial constante, el trabajo
sale \SI{32.7}{\giga\joule} y el combustible \SI{2547}{\kilogram}, que es
exactamente el resultado de la forma lineal: la aproximación lineal
\emph{es} suponer que el avión no adelgaza. Si se usa la media logarítmica de
los pesos, \SI{8828}{\kilogram}, el trabajo sale \SI{28.9}{\giga\joule} y el
combustible \SI{2248}{\kilogram}, que reproduce la
ecuación~\eqref{eq:resultado} con cinco cifras. La media logarítmica no es una
elección estética: es la media que la integral del logaritmo impone.

\textbf{Comprobación 2: potencias y flujo de combustible.} La potencia
propulsiva media es \(\overline{D}\,V = \SI{5774}{\newton} \times
\SI{300}{\meter\per\second} = \SI{1.73}{\mega\watt}\), y la potencia química
que hay que meter por el otro extremo es \SI{5.77}{\mega\watt}. Dividida por
la densidad energética, eso son \SI{0.135}{\kilogram\per\second}, es decir
\SI{486}{\kilogram\per\hour}, que multiplicados por las \SI{4.63}{\hour} de
crucero devuelven los \SI{2248}{\kilogram}. El consumo pasa de
\SI{550}{\kilogram\per\hour} al principio a \SI{426}{\kilogram\per\hour} al
final: un 23 por ciento menos, sólo porque el avión pesa menos.

\textbf{Comprobación 3: la resistencia.} La ecuación~\eqref{eq:endurance} da
\(E = R/V = \SI{5.0e6}{\meter}/\SI{300}{\meter\per\second} = \SI{16667}{\second}
= \SI{4.63}{\hour}\), es decir 4 horas y 38 minutos. Aquí es donde entra la
velocidad, que en el alcance se había cancelado. Es una comprobación débil (no
usa la ecuación de Bréguet, sólo la definición de velocidad) pero es la que
contesta la mitad \enquote{de resistencia} de la pregunta.

\textbf{Comprobación 4: contra un avión que existe.} Un avión de transporte con
peso máximo al despegue de \SI{9752}{\kilogram}, que es la masa de este
problema con un 2 por ciento de diferencia, lleva \SI{2750}{\kilogram} de
combustible y tiene un alcance certificado de \SI{3797}{\kilo\meter} con cuatro
pasajeros \parencite{learjet75,aopalearjet}. Eso son
\SI{0.724}{\kilogram\per\kilo\meter}, frente a los
\SI{0.450}{\kilogram\per\kilo\meter} de mi estimación: la estimación es
\textbf{optimista en un factor 1,6}. Invirtiendo la ecuación de Bréguet con los
datos del avión real se obtiene \(\eta_o (L/D) = \num{2.63}\), frente a
\(0{,}30\times 15 = \num{4.50}\) que supone el enunciado. Esa es la medida
limpia de cuánto de bueno es el avión hipotético del enunciado comparado con
uno de verdad, y la respuesta es que es un 71 por ciento mejor de lo que existe.

\begin{table}[htbp]
\small
\caption{Las cuatro comprobaciones y lo que cada una añade. Cálculo propio.}
\label{cuadro:comprobaciones}
\begin{tabularx}{\textwidth}{@{}p{4.0cm}rXl@{}}
\toprule
\textbf{Comprobación} & \textbf{Resultado} & \textbf{Qué verifica} &
\textbf{Veredicto} \\
\midrule
Trabajo contra el arrastre con peso medio logarítmico
& \SI{2248}{\kilogram}
& que la integral está bien hecha
& \textcolor{cVerde}{\textbf{coincide}} \\
Trabajo contra el arrastre con peso inicial constante
& \SI{2547}{\kilogram}
& que la forma lineal equivale a ignorar el adelgazamiento
& \textcolor{cVerde}{\textbf{explica el} \boldmath\SI{13}{\percent}} \\
Potencia y flujo de combustible medio
& \SI{486}{\kilogram\per\hour}
& que las magnitudes intermedias son físicamente razonables
& \textcolor{cVerde}{\textbf{coincide}} \\
Consumo por kilómetro de un reactor real de \SI{9.8}{\tonne}
& \SI{0.724}{\kilogram\per\kilo\meter}
& que el orden de magnitud es el del mundo real
& \textcolor{cAviso}{\textbf{estimación} \boldmath\num{1.6}\textbf{ veces
  baja}} \\
\bottomrule
\end{tabularx}
\end{table}

\section{Para qué sirven la velocidad y la densidad del aire}
\label{sec:coherencia}

El enunciado da dos datos que \textbf{no entran en la ecuación del alcance}: la
velocidad de \SI{300}{\meter\per\second} y la densidad del aire de
\SI{0.6125}{\kilogram\per\meter\cubed}. La velocidad se cancela al pasar de la
resistencia al alcance, y la densidad no aparece en ninguna de las dos
ecuaciones. No son datos superfluos: son los datos que permiten comprobar si el
resto del enunciado describe un avión posible. Y la respuesta, con números, es
que no.

\subsection{Primer indicio: el número de Mach}

A \SI{6000}{\meter} en atmósfera estándar la temperatura es
\SI{249.15}{\kelvin} \parencite{iso2533}, de donde la velocidad del sonido es
\begin{equation}
a = \sqrt{\gamma R_{\text{aire}} T}
  = \sqrt{1{,}4 \times \SI{287.05}{\joule\per\kilogram\per\kelvin}
          \times \SI{249.15}{\kelvin}}
  = \SI{316.4}{\meter\per\second},
\end{equation}
y por tanto el número de Mach del enunciado es
\begin{equation}
M = \frac{V}{a} = \frac{300}{316{,}4} = \num{0.948}.
\end{equation}
Eso es un problema. Los aviones de transporte cruzan entre Mach 0,78 y Mach
0,85 precisamente porque por encima de Mach 0,90 aparece la
\textbf{divergencia de arrastre}: se forman ondas de choque locales sobre el
extradós del ala y el coeficiente de arrastre se dispara
\parencite{anderson1999}. Mach 0,948 está por encima del Mach de divergencia de
cualquier ala subsónica conocida. Un \(L/D\) de 15 en esa condición no es
optimista: es imposible.

De paso, la densidad que da el enunciado también merece un comentario. La
atmósfera estándar a \SI{6000}{\meter} tiene
\(\rho = \SI{0.660}{\kilogram\per\meter\cubed}\), que es el 53,8 por ciento del
valor a nivel del mar y no el 50 por ciento \parencite{iso2533}. La
aproximación del enunciado es buena (un 7 por ciento de diferencia) y no cambia
nada, pero conviene saber que \(\rho_0/2\) corresponde en realidad a
\SI{6600}{\meter}.

\subsection{Segundo indicio: la sustentación no cierra}

Con la densidad y la velocidad del enunciado, la presión dinámica es
\begin{equation}
q = \tfrac{1}{2}\rho V^{2}
  = \tfrac{1}{2}\times\SI{0.6125}{\kilogram\per\meter\cubed}
    \times(\SI{300}{\meter\per\second})^{2}
  = \SI{27562}{\pascal},
\end{equation}
casi tres veces la presión dinámica de crucero de un avión de línea
(\SI{9600}{\pascal} para un reactor de pasillo único a \SI{11000}{\meter} y
Mach 0,78). De la ecuación de sustentación \(W = q\,S\,C_L\) sale entonces
\begin{equation}
S\,C_L = \frac{W}{q} = \frac{\SI{98100}{\newton}}{\SI{27562}{\pascal}}
       = \SI{3.56}{\meter\squared},
\end{equation}
y aquí se abre la tenaza. El Cuadro~\ref{cuadro:coherencia} y la
Figura~\ref{fig:coherencia} la desarrollan:

\begin{itemize}
\item Si el ala es de tamaño normal para un avión de \SI{10}{\tonne}, del orden
  de \SI{29}{\meter\squared} \parencite{learjet75}, entonces
  \(C_L = 3{,}56/29 = \num{0.123}\). Para que \(L/D = 15\) con ese
  \(C_L\) haría falta \(C_D = 0{,}123/15 = \num{0.0082}\), que es
  \textbf{menos de la mitad del arrastre parásito} de cualquier avión de
  transporte, cuyo \(C_{D0}\) mínimo ronda \num{0.020} a \num{0.025}
  \parencite{anderson1999,raymer2018}. Con \(C_{D0} = \num{0.022}\) el
  \(L/D\) realmente alcanzable en esa condición es \num{5.6}.
\item Si en cambio se impone un \(C_L\) razonable de crucero, del orden de
  \num{0.40}, entonces el ala tendría que medir \SI{8.9}{\meter\squared}, lo
  que da una carga alar de \SI{1124}{\kilogram\per\meter\squared}. Eso es más
  que la de un caza de combate.
\end{itemize}

Las dos salidas son inadmisibles, y la conclusión es que \textbf{los datos del
enunciado no describen ningún avión construible}: describen un avión de
\SI{10}{\tonne} con la eficiencia aerodinámica de un planeador volando a la
velocidad de un caza a una altitud de avión de hélice. Esto no invalida el
ejercicio, porque el ejercicio es sobre la ecuación y no sobre el avión, pero sí
es una de las respuestas a la pregunta de por qué el cálculo es complicado, y de
las más importantes: en un problema real \textbf{no se pueden elegir los datos
de entrada por separado}, porque están acoplados por la física.

\begin{table}[htbp]
\small
\caption{Examen de coherencia entre los datos del enunciado y la aerodinámica.
Cálculo propio; valores de referencia de reactor de pasillo único a
\SI{11000}{\meter} y Mach 0,78, y coeficiente de arrastre parásito de
\textcite{anderson1999} y \textcite{raymer2018}.}
\label{cuadro:coherencia}
\begin{tabularx}{\textwidth}{@{}Xrrl@{}}
\toprule
\textbf{Magnitud} & \textbf{Enunciado} & \textbf{Avión real} &
\textbf{Diagnóstico} \\
\midrule
Altitud (\si{\meter})                          & \num{6000}  & \num{11000} &
  baja para crucero \\
Número de Mach                                 & \num{0.948} & \num{0.78}  &
  \textcolor{cAviso}{transónico} \\
Presión dinámica \(q\) (\si{\pascal})          & \num{27562} & \num{9600}  &
  \(\times 2{,}9\) \\
Coeficiente de sustentación \(C_L\)            & \num{0.123} & \num{0.540}  &
  \textcolor{cAviso}{muy bajo} \\
\(C_D\) necesario para \(L/D = 15\)            & \num{0.0082} & \num{0.032} &
  \textcolor{cAviso}{por debajo del parásito} \\
\(L/D\) realmente alcanzable                   & \num{5.6}   & \num{17}    &
  \textcolor{cAviso}{un tercio de lo supuesto} \\
\midrule
Combustible con \(L/D = 15\) (\si{\kilogram})  & \num{2248}  &             &
  respuesta del ejercicio \\
Combustible con \(L/D = 5{,}6\) (\si{\kilogram}) & \num{4958} &            &
  \textcolor{cAviso}{el 50 \% del peso al despegue} \\
Alcance real con \SI{2750}{\kilogram} y \(L/D = 5{,}6\) (\si{\kilo\meter})
  & \num{2419} & \num{3797} & \textcolor{cAviso}{no llega} \\
\bottomrule
\end{tabularx}
\end{table}

\begin{figure}[htbp]
\centering
\begin{tikzpicture}

% ---------------- Panel A: la condición de vuelo ----------------
\node[anchor=south west, font=\small\bfseries] at (-0.5,4.30)
     {Panel A. La condición de vuelo del enunciado frente a la de un reactor
      de línea};

% presión dinámica
\foreach \v in {10,20,30} {
  \draw[cNeutro!28] (0,{\v*0.125}) -- (5.6,{\v*0.125});
  \node[left, font=\scriptsize, text=cNeutro] at (-0.06,{\v*0.125}) {\v};
}
\draw[cNeutro!70] (0,0) -- (5.6,0);
\fill[cDato!75] (0.55,0) rectangle (1.75,{9.6*0.125});
\node[above, font=\scriptsize\bfseries, text=cDato] at (1.15,{9.6*0.125})
     {\num{9.6}};
\fill[cAviso!80] (2.75,0) rectangle (3.95,{27.562*0.125});
\node[above, font=\scriptsize\bfseries, text=cAviso] at (3.35,{27.562*0.125})
     {\num{27.6}};
\node[below, font=\scriptsize, text width=2.2cm, align=center] at (1.15,-0.10)
     {reactor de línea\\M 0,78 a \SI{11}{\kilo\meter}};
\node[below, font=\scriptsize, text width=2.2cm, align=center] at (3.35,-0.10)
     {\textbf{enunciado}\\M 0,95 a \SI{6}{\kilo\meter}};
\node[font=\scriptsize, text=cNeutro, rotate=90] at (-0.80,2.1)
     {presión dinámica $q$ (kPa)};

% coeficiente de sustentación
\begin{scope}[shift={(7.6,0)}]
\foreach \v/\t in {0.2/{0,2}, 0.4/{0,4}, 0.6/{0,6}} {
  \draw[cNeutro!28] (0,{\v*6.2}) -- (5.6,{\v*6.2});
  \node[left, font=\scriptsize, text=cNeutro] at (-0.06,{\v*6.2}) {\t};
}
\draw[cNeutro!70] (0,0) -- (5.6,0);
\fill[cDato!75] (0.55,0) rectangle (1.75,{0.540*6.2});
\node[above, font=\scriptsize\bfseries, text=cDato] at (1.15,{0.540*6.2})
     {\num{0.54}};
\fill[cAviso!80] (2.75,0) rectangle (3.95,{0.123*6.2});
\node[above, font=\scriptsize\bfseries, text=cAviso] at (3.35,{0.123*6.2})
     {\num{0.123}};
\node[below, font=\scriptsize, text width=2.2cm, align=center] at (1.15,-0.10)
     {reactor de línea\\$S = \SI{123}{\meter\squared}$};
\node[below, font=\scriptsize, text width=2.2cm, align=center] at (3.35,-0.10)
     {\textbf{enunciado}\\$S = \SI{29}{\meter\squared}$};
\node[font=\scriptsize, text=cNeutro, rotate=90] at (-0.75,2.1)
     {coeficiente de sustentación $C_L$};
\node[right, font=\scriptsize, text=cAviso, align=left, text width=3.0cm]
     at (4.15,2.5)
     {volar tan rápido y tan bajo obliga a un $C_L$ cuatro veces menor
      que el de crucero};
\end{scope}

% ---------------- Panel B: la consecuencia ----------------
\begin{scope}[shift={(0,-7.35)}]
\node[anchor=south west, font=\small\bfseries] at (-0.9,5.55)
     {Panel B. Consecuencia: el $L/D$ supuesto y el combustible que hace falta
      de verdad};

% L/D
\foreach \v in {5,10,15,20} {
  \draw[cNeutro!28] (0,{\v*0.25}) -- (5.6,{\v*0.25});
  \node[left, font=\scriptsize, text=cNeutro] at (-0.06,{\v*0.25}) {\v};
}
\draw[cNeutro!70] (0,0) -- (5.6,0);
\fill[cDato!75] (0.35,0) rectangle (1.35,{17*0.25});
\node[above, font=\scriptsize\bfseries, text=cDato] at (0.85,{17*0.25}) {17};
\fill[cVerde!80] (2.05,0) rectangle (3.05,{15*0.25});
\node[above, font=\scriptsize\bfseries, text=cVerde] at (2.55,{15*0.25}) {15};
\fill[cAviso!80] (3.75,0) rectangle (4.75,{5.6*0.25});
\node[above, font=\scriptsize\bfseries, text=cAviso] at (4.25,{5.6*0.25})
     {\num{5.6}};
\node[below, font=\scriptsize, text width=1.75cm, align=center] at (0.85,-0.10)
     {reactor de línea (real)};
\node[below, font=\scriptsize, text width=1.75cm, align=center] at (2.55,-0.10)
     {\textbf{supuesto en el enunciado}};
\node[below, font=\scriptsize, text width=1.75cm, align=center] at (4.25,-0.10)
     {alcanzable a M 0,95 con $C_{D0}=0{,}022$};
\node[font=\scriptsize, text=cNeutro, rotate=90] at (-0.70,2.1)
     {relación $L/D$};

% combustible
\begin{scope}[shift={(7.6,0)}]
\foreach \v in {1000,2000,3000,4000,5000} {
  \draw[cNeutro!28] (0,{\v*0.00098}) -- (5.6,{\v*0.00098});
  \node[left, font=\scriptsize, text=cNeutro] at (-0.06,{\v*0.00098})
       {\num{\v}};
}
\draw[cNeutro!70] (0,0) -- (5.6,0);
\fill[cVerde!80] (0.75,0) rectangle (2.05,{2248*0.00098});
\node[above, font=\scriptsize\bfseries, text=cVerde] at (1.40,{2248*0.00098})
     {\num{2248}};
\fill[cAviso!80] (3.15,0) rectangle (4.45,{4958*0.00098});
\node[above, font=\scriptsize\bfseries, text=cAviso] at (3.80,{4958*0.00098})
     {\num{4958}};
\draw[cDato, very thick, dashed] (0,{2750*0.00098}) -- (5.6,{2750*0.00098});
\node[above right, font=\scriptsize, text=cDato, align=left]
     at (0.10,{2750*0.00098})
     {depósito: \SI{2750}{\kilogram}};
\node[below, font=\scriptsize, text width=2.1cm, align=center] at (1.40,-0.10)
     {con $L/D = 15$\\(respuesta del ejercicio)};
\node[below, font=\scriptsize, text width=2.1cm, align=center] at (3.80,-0.10)
     {con $L/D = 5{,}6$\\(el \SI{50}{\percent} del avión)};
\node[font=\scriptsize, text=cNeutro, rotate=90] at (-1.32,2.4)
     {queroseno (kg)};
\end{scope}
\end{scope}

\end{tikzpicture}
\caption{Por qué la velocidad y la densidad, que no entran en la ecuación,
importan. Panel A: a \SI{300}{\meter\per\second} y \SI{6000}{\meter} la presión
dinámica es 2,9 veces la de crucero de un reactor de línea, lo que fuerza un
coeficiente de sustentación cuatro veces menor. Panel B: con ese \(C_L\) y un
arrastre parásito realista, el \(L/D\) alcanzable es \num{5.6} y no 15, y el
combustible necesario pasa de \num{2248} a \SI{4958}{\kilogram}, por encima de
la capacidad de depósito de cualquier avión de esta masa. Los datos del
enunciado son aritméticamente consistentes y físicamente incompatibles entre
sí. Elaboración propia; \(C_{D0}\) de \textcite{anderson1999} y
\textcite{raymer2018}, atmósfera de \textcite{iso2533}.}
\label{fig:coherencia}
\end{figure}

\section{Del crucero ideal al depósito real}
\label{sec:cascada}

El enunciado autoriza a pasar por alto el ascenso y el descenso, y hago uso de
esa licencia para el resultado principal. Pero conviene saber cuánto vale la
licencia, porque es la diferencia entre una estimación y un plan de vuelo. La
Figura~\ref{fig:cascada} pone la cuenta completa.

\begin{figure}[htbp]
\centering
\begin{tikzpicture}

\node[anchor=south west, font=\small\bfseries] at (-0.9,5.80)
     {De la respuesta del ejercicio al combustible que habría que embarcar de
      verdad (kg)};

% rejilla
\foreach \v in {500,1000,1500,2000,2500,3000} {
  \draw[cNeutro!26] (0,{\cWy{\v}}) -- (10.15,{\cWy{\v}});
  \node[left, font=\scriptsize, text=cNeutro] at (-0.06,{\cWy{\v}})
       {\num{\v}};
}
\draw[cNeutro!70] (0,0) -- (10.15,0);
\draw[cNeutro!70] (0,0) -- (0,5.55);

% capacidad de depósito de un reactor real de 10 t
\draw[cAviso, very thick, dashed] (0,{\cWy{2750}}) -- (10.15,{\cWy{2750}});
\node[right, font=\scriptsize, text=cAviso, align=left, text width=2.75cm]
     at (10.35,{\cWy{2620}})
     {capacidad de depósito de un reactor real de \SI{9752}{\kilogram}:
      \SI{2750}{\kilogram}};

% líneas de enlace entre tramos
\foreach \x/\y in {1.44/2248, 2.84/2329, 4.24/2441, 5.64/2607, 7.04/2807,
                   8.44/2847} {
  \draw[cNeutro!60, dotted] (\x,{\cWy{\y}}) -- ({\x+0.26},{\cWy{\y}});
}

\tramo{0.30}{0}{2248}{cCalc}{crucero ideal (Bréguet)}{\num{2248}}{-0.14}
\tramo{1.70}{2248}{2329}{cNaranja}{ascenso y aceleración}{$+81$}{-0.86}
\tramo{3.10}{2329}{2441}{cNaranja}{contingencia del
  \SI{5}{\percent}}{$+112$}{-0.14}
\tramo{4.50}{2441}{2607}{cNaranja}{alternativo a
  \SI{370}{\kilo\meter}}{$+166$}{-0.86}
\tramo{5.90}{2607}{2807}{cNaranja}{reserva final de
  \SI{30}{\minute}}{$+200$}{-0.14}
\tramo{7.30}{2807}{2847}{cNaranja}{rodaje}{$+40$}{-0.86}
\tramo{8.86}{0}{2847}{cAviso}{\textbf{total a embarcar}}{\num{2847}}{-0.14}

\node[anchor=north west, font=\scriptsize, align=left, text width=13.0cm,
      text=cNeutro!85!black] at (-0.9,-1.70)
     {La licencia que da el enunciado (\enquote{puede pasar por alto el ascenso
      y el descenso}) vale \SI{81}{\kilogram}, un \SI{3.6}{\percent}. Lo que de
      verdad falta es la normativa de reservas: \SI{518}{\kilogram}, un
      \SI{23}{\percent}, y no es negociable.};

\end{tikzpicture}
\caption{La respuesta del ejercicio es la primera barra; las cinco siguientes
son lo que la normativa y la física añaden. El incremento por ascenso es el
mínimo termodinámico, calculado como la energía potencial más la cinética que
hay que darle al avión
(\(mgh + \tfrac{1}{2}mV^{2} = \num{589} + \num{450} = \SI{1039}{\mega\joule}\),
que a \(\eta_o h\) por kilogramo son \SI{81}{\kilogram}); en la práctica es
mayor, porque durante el ascenso ni el \(L/D\) ni el \(\eta_o\) valen lo que
valen en crucero. Las reservas siguen la política de combustible del Reglamento
(UE) 965/2012 \parencite{easa965}: contingencia del 5 por ciento del trayecto,
combustible hasta un aeropuerto alternativo a \num{200} millas náuticas,
reserva final de 30 minutos en espera y rodaje. Elaboración propia. El total,
\SI{2847}{\kilogram}, supera la capacidad de depósito del avión real de masa
comparable, y eso con el \(L/D\) optimista del enunciado.}
\label{fig:cascada}
\end{figure}

Los dos números que salen de la cascada son estos. La licencia del enunciado
(ignorar ascenso y descenso) vale \textbf{\SI{81}{\kilogram}, un 3,6 por
ciento}, y por tanto está bien concedida: no cambia la respuesta. Lo que sí la
cambia es la normativa: las reservas obligatorias suman
\textbf{\SI{518}{\kilogram}, un 23 por ciento} sobre el combustible de
trayecto, y no son opcionales. Un avión que despachara con
\SI{2248}{\kilogram} para este vuelo no despegaría.

Y el resultado que cierra el círculo: el total, \SI{2847}{\kilogram}, es
\textbf{superior a la capacidad de depósito} del reactor real de masa
equivalente, \SI{2750}{\kilogram} \parencite{learjet75}. Es decir que incluso
concediéndole al avión del enunciado su \(L/D\) imposible de 15, el vuelo
Boston--Los Ángeles sin escalas queda fuera de alcance para un avión de
\SI{10}{\tonne}. Con el \(L/D\) real de la sección~\ref{sec:coherencia} queda
fuera por más del doble. La ecuación de Bréguet, bien usada, no sólo da un
número: dice que el avión del enunciado no puede hacer el viaje del enunciado.

\section{Análisis de sensibilidad}

Como los datos de entrada son estimaciones y no medidas, lo que importa no es
la cifra sino su intervalo. La Figura~\ref{fig:tornado} y el
Cuadro~\ref{cuadro:sensibilidad} lo dan.

\begin{figure}[htbp]
\centering
\begin{tikzpicture}

\node[anchor=south west, font=\small\bfseries] at (-3.7,5.10)
     {Efecto de cada dato de entrada sobre el queroseno estimado (kg)};

% rejilla
\foreach \v in {1800,2000,2200,2400,2600,2800,3000} {
  \draw[cNeutro!26] ({\tSx{\v}},-0.10) -- ({\tSx{\v}},4.55);
  \node[below, font=\scriptsize, text=cNeutro] at ({\tSx{\v}},-0.12)
       {\num{\v}};
}
% línea del caso base
\draw[cDato, very thick] ({\tSx{2248}},-0.10) -- ({\tSx{2248}},4.62);
\node[above, font=\scriptsize\bfseries, text=cDato] at ({\tSx{2248}},4.60)
     {caso base \num{2248}};

\tornado{4.0}{1739}{2633}{Eficiencia total $\eta_o$\\(de \num{0.25} a
  \num{0.40})}{\num{0.40}}{\num{0.25}}
\tornado{3.1}{1912}{2726}{Relación $L/D$\\(de 12 a 18)}{18}{12}
\tornado{2.2}{2248}{2900}{Lectura del enunciado\\(peso total o peso sin
  combustible)}{A}{B}
\tornado{1.3}{1923}{2443}{Distancia $R$\\(de \num{4193} a
  \SI{5500}{\kilo\meter})}{\num{4193}}{\num{5500}}
\tornado{0.4}{2094}{2262}{Densidad energética $h$\\(de \num{42.5} a
  \SI{46.4}{\mega\joule\per\kilogram})}{\num{46.4}}{\num{42.5}}

\node[font=\scriptsize, align=left, text width=4.4cm,
      text=cNeutro!85!black] at (10.3,0.80)
     {El orden importa: los dos datos que más pesan son los dos que el
      enunciado regala hechos.};

\end{tikzpicture}
\caption{Diagrama de tornado del combustible estimado. Cada barra recorre el
intervalo plausible de un dato de entrada manteniendo los demás en su valor
base. Los dos primeros, la eficiencia total y la relación sustentación/arrastre,
son los que dominan y son precisamente los dos que en un problema real hay que
estimar con un modelo del motor y un modelo del ala; el enunciado los da por
supuestos, y ahí es donde se esconde la dificultad del cálculo. El tercero es
la ambigüedad de una frase del propio enunciado. Elaboración propia.}
\label{fig:tornado}
\end{figure}

\begin{table}[htbp]
\small
\caption{Sensibilidad del resultado a cada dato de entrada, lectura A. Cada
fila cambia un solo dato. Cálculo propio.}
\label{cuadro:sensibilidad}
\begin{tabularx}{\textwidth}{@{}Xrrl@{}}
\toprule
\textbf{Variante} & \textbf{Queroseno} & \textbf{Cambio} &
\textbf{Comentario} \\
 & \textbf{(\si{\kilogram})} & & \\
\midrule
\textbf{Caso base} (\(h=42{,}8\); \(\eta_o=0{,}30\); \(L/D=15\);
  \(R=\SI{5000}{\kilo\meter}\))
  & \textbf{\num{2248}} & & respuesta del ejercicio \\
\midrule
\(h = \SI{42.5}{\mega\joule\per\kilogram}\) (valor del módulo)
  & \num{2262} & \num[retain-explicit-plus]{+0.6}\,\% & indiferente \\
\(h = \SI{43.0}{\mega\joule\per\kilogram}\) (Jet A-1 típico)
  & \num{2239} & \num{-0.4}\,\% & indiferente \\
\(h = \SI{46.4}{\mega\joule\per\kilogram}\) (poder calorífico superior)
  & \num{2094} & \num{-6.9}\,\% & \textcolor{cAviso}{error conceptual} \\
\midrule
\(\eta_o = 0{,}25\) (turbofán antiguo)
  & \num{2633} & \num[retain-explicit-plus]{+17.1}\,\% & \\
\(\eta_o = 0{,}35\) (turbofán moderno)
  & \num{1961} & \num{-12.8}\,\% & \\
\(\eta_o = 0{,}40\) (mejor disponible)
  & \num{1739} & \num{-22.7}\,\% & \\
\midrule
\(L/D = 12\) (reactor de negocios real)
  & \num{2726} & \num[retain-explicit-plus]{+21.3}\,\% & \\
\(L/D = 18\) (reactor de línea moderno)
  & \num{1912} & \num{-14.9}\,\% & \\
\(L/D = 5{,}6\) (alcanzable a Mach 0,95)
  & \num{4958} & \num[retain-explicit-plus]{+120.5}\,\% &
  \textcolor{cAviso}{ver Cuadro~\ref{cuadro:coherencia}} \\
\midrule
\(R = \SI{4193}{\kilo\meter}\) (distancia real BOS--LAX)
  & \num{1923} & \num{-14.9}\,\% &
  \textcolor{cAviso}{el enunciado redondea un 19 \% al alza} \\
\(R = \SI{5500}{\kilo\meter}\) (desvío y viento en contra del 10 \%)
  & \num{2443} & \num[retain-explicit-plus]{+8.7}\,\% & \\
\midrule
Lectura B (\SI{10}{\tonne} sin combustible)
  & \num{2900} & \num[retain-explicit-plus]{+29.0}\,\% & ambigüedad del
  enunciado \\
Forma lineal de la ecuación
  & \num{2547} & \num[retain-explicit-plus]{+13.3}\,\% & aproximación de primer
  orden \\
Con reservas normativas (Figura~\ref{fig:cascada})
  & \num{2847} & \num[retain-explicit-plus]{+26.6}\,\% & obligatorio en
  operación real \\
\bottomrule
\end{tabularx}
\end{table}

Dos conclusiones del cuadro. La primera: \textbf{el dato que el enunciado nos
manda buscar es el que menos importa}. Cambiar la densidad energética del
queroseno dentro de todo su rango de especificación mueve el resultado menos
del 1 por ciento, mientras que cambiar la eficiencia total en un punto
razonable lo mueve 20 veces más. La segunda: \textbf{la distancia del enunciado
no es la distancia}. Boston y Los Ángeles están a \SI{4193}{\kilo\meter} de
círculo máximo; los \SI{5000}{\kilo\meter} del enunciado son un 19 por ciento
más, que probablemente sea una concesión razonable al encaminamiento real y al
viento, pero conviene declararla, porque por sí sola vale más que toda la
incertidumbre del combustible.

\section{De queroseno a huella de carbono por pasajero}

La actividad se titula \enquote{estima la huella de carbono por pasajero de un
vuelo}, así que cierro el círculo. La combustión completa de un kilogramo de
queroseno produce \SI{3.16}{\kilogram} de \ce{CO2}, cifra que no depende del
motor sino de la estequiometría del hidrocarburo \parencite{icaocalc}. Con los
\SI{2248}{\kilogram} de la respuesta:
\begin{equation}
m_{\ce{CO2}} = \SI{2248}{\kilogram}\times\num{3.16}
             = \SI{7105}{\kilogram} \approx \SI{7.1}{\tonne}.
\end{equation}

El Cuadro~\ref{cuadro:huella} reparte esa cifra entre pasajeros. El reparto es
la parte interesante, porque es donde se ve que \textbf{la huella por pasajero
no es una propiedad del avión sino del modelo de negocio}.

\begin{table}[htbp]
\small
\caption{La misma tonelada de queroseno, repartida entre distintos números de
pasajeros. Vuelo de \SI{5000}{\kilo\meter}, sólo crucero, sin reservas y con
factor de ocupación del 100 por ciento. Cálculo propio; factor de emisión de
\textcite{icaocalc}, valores de referencia de \textcite{defra2025}.}
\label{cuadro:huella}
\begin{tabularx}{\textwidth}{@{}Xrrrr@{}}
\toprule
\textbf{Configuración} & \textbf{Pasajeros} & \textbf{\ce{CO2} por} &
\textbf{\ce{CO2} por} & \textbf{Queroseno} \\
 & & \textbf{pasajero} & \textbf{\si{\pkm}} & \textbf{(\si{\litro\per\pkm})} \\
 & & \textbf{(\si{\kilogram})} & \textbf{(\si{\kilogram})} & \\
\midrule
Reactor privado, cabina de lujo & 4 & \num{1776} & \num{0.355} & \num{0.140} \\
Reactor privado, uso habitual   & 6 & \num{1184} & \num{0.237} & \num{0.093} \\
Reactor de negocios, cabina llena & 8 & \num{888} & \num{0.178} & \num{0.070} \\
Configuración máxima de la clase & 9 & \num{789} & \num{0.158} & \num{0.062} \\
\midrule
\emph{Referencia}: reactor de pasillo único, \num{180} plazas
  & \num{180} & \num{243} & \num{0.049} & \num{0.019} \\
\emph{Referencia}: factor oficial británico, vuelo de largo radio
  & & & \num{0.117} & \\
\emph{Referencia}: factor oficial británico, vuelo de radio corto
  & & & \num{0.150} & \\
\bottomrule
\end{tabularx}
\end{table}

La lectura es incómoda y es correcta: \textbf{el mismo vuelo, en el mismo
avión, tiene una huella por pasajero que varía en un factor 2,25 según cuánta
gente vaya dentro}, y comparado con un reactor de pasillo único lleno la
diferencia llega a un factor 3,6 a favor del avión grande. La eficiencia
aeronáutica del avión pequeño es peor, pero lo que domina el resultado no es la
aerodinámica: es el denominador.

Falta una advertencia. Estos \SI{7.1}{\tonne} son sólo \ce{CO2}. Si se
contabiliza el forzamiento radiativo completo, incluidas las estelas de
condensación, los óxidos de nitrógeno y el vapor de agua, el efecto climático
total de la aviación es del orden de tres veces el del \ce{CO2} por sí solo
\parencite{lee2021}, lo que llevaría este vuelo a unas \SI{21}{\tonne} de
\ce{CO2} equivalente. Ese multiplicador depende de la altitud, de la latitud,
de la hora del día y de la humedad de la capa de aire que se atraviesa, y por
tanto \textbf{la huella real de este vuelo no está determinada por el
queroseno que se quema}. Es la última y quizá la más incómoda de las razones
por las que el cálculo es complicado.

\section{¿Por qué es un cálculo complicado?}
\label{sec:complicado}

La ecuación se resuelve en cinco líneas y la aritmética es la de una
calculadora de bolsillo. Lo complicado no es la ecuación: es que
\textbf{ninguno de sus cuatro factores es un dato, todos son resultados de otro
problema}, y que el propio problema está definido de forma implícita. Doy diez
razones, cada una con su magnitud medida en este ejercicio, y las resumo en el
Cuadro~\ref{cuadro:porque}.

\subsection{Razón 1: el combustible pesa, y hay que llevarlo}

Es la dificultad estructural del problema y está en la
ecuación~\eqref{eq:edo}: el peso aparece a la vez como lo que hay que sostener
y como lo que se está consumiendo. De ahí el logaritmo, y de ahí que el
problema sea \textbf{implícito}: no se puede calcular el combustible sin saber
el peso, y no se puede saber el peso sin conocer el combustible.

La consecuencia práctica se mide con la derivada de la
ecuación~\eqref{eq:breguet} respecto a la carga de pago:
\begin{equation}
\frac{d\,m_{\text{comb}}}{d\,m_{\text{pago}}} = e^{x}-1 = \num{0.290}.
\end{equation}
Cada kilogramo de carga extra en este vuelo obliga a embarcar
\SI{0.29}{\kilogram} más de queroseno, y por tanto el peso al despegue sube
\SI{1.29}{\kilogram}. Es el \textbf{efecto bola de nieve} del diseño
aeronáutico. En un vuelo de \SI{15000}{\kilo\meter} el mismo factor vale
\num{1.15}, es decir que un kilogramo de carga cuesta más de un kilogramo de
combustible, y en un avión de diseño el efecto se realimenta también a través
del ala y del tren de aterrizaje, que hay que agrandar. Por eso el diseño
preliminar de un avión es un bucle iterativo y no una fórmula
\parencite{raymer2018}.

\subsection{Razón 2: la ecuación sólo vale en la condición en la que no se
vuela}

La ecuación~\eqref{eq:breguet} se obtuvo integrando con \(\eta_o\), \(L/D\) y
\(V\) constantes. Un vuelo real no tiene ninguna de las tres constantes. A
medida que el avión adelgaza, el \(C_L\) necesario baja, y con él el punto de
la polar en el que se vuela y por tanto el \(L/D\); la práctica habitual es
subir de altitud en escalones (\emph{step climb}) precisamente para mantener el
\(L/D\) cerca de su óptimo, lo que a su vez cambia la densidad, el número de
Mach y el rendimiento del motor. La integral exacta requiere la polar del ala y
el mapa del motor, y entonces ya no es una integral cerrada: es una simulación
paso a paso.

\subsection{Razón 3: hay fases de vuelo que la ecuación no describe}

El enunciado autoriza a ignorar el ascenso y el descenso, y en este caso la
autorización es benigna. El coste energético mínimo de subir a
\SI{6000}{\meter} y acelerar a \SI{300}{\meter\per\second} es
\begin{equation}
m\,g\,h + \tfrac{1}{2}m V^{2}
 = \SI{589}{\mega\joule} + \SI{450}{\mega\joule} = \SI{1039}{\mega\joule},
\end{equation}
que a \(\eta_o h = \SI{12.84}{\mega\joule\per\kilogram}\) de trabajo útil por
kilogramo de queroseno son \SI{81}{\kilogram}, un 3,6 por ciento del total.
Pero ese es el mínimo termodinámico: durante el ascenso el avión vuela con
\(C_L\) alto, arrastre inducido alto y motor fuera de su punto óptimo, y el
consumo real de la fase de ascenso de un vuelo de este tipo es del orden del
doble. En vuelos cortos la fase de ascenso puede ser la mitad del consumo total,
y ahí la ecuación de Bréguet directamente no sirve.

\subsection{Razón 4: la normativa manda más que la física}

Ningún avión despega con el combustible que necesita. Despega con el
combustible que necesita más la contingencia, más el alternativo, más la
reserva final, más el rodaje. En este vuelo eso son \SI{518}{\kilogram}
adicionales, un \textbf{23 por ciento} \parencite{easa965}, y ese 23 por ciento
tiene que estar ahí aunque el vuelo transcurra sin novedad. La
Figura~\ref{fig:cascada} lo detalla. Es un caso claro de una restricción que no
está en ninguna ecuación de la mecánica del vuelo y que sin embargo determina si
el avión puede hacer el viaje.

\subsection{Razón 5: los dos factores que dominan son los dos más difíciles de
conocer}

El Cuadro~\ref{cuadro:sensibilidad} lo cuantifica: \(\eta_o\) y \(L/D\) mueven
el resultado un 40 y un 36 por ciento respectivamente en su rango plausible,
mientras que la densidad energética lo mueve un 1 por ciento. Y \(\eta_o\) y
\(L/D\) no son datos tabulados: son funciones del número de Mach, de la
altitud, del peso, del régimen del motor, de la temperatura del aire y del
estado de la superficie del avión. Obtenerlos con precisión exige un modelo
aerodinámico y un modelo de motor, es decir exige haber resuelto ya el problema
de diseño. Se acaba en la situación paradójica de que \textbf{la parte fácil del
cálculo es la ecuación y la parte difícil son los datos que hay que meterle}.

\subsection{Razón 6: la densidad energética no es un número, es un intervalo}

El poder calorífico inferior del Jet A-1 varía entre \num{42.8} y
\SI{43.2}{\mega\joule\per\kilogram} según el lote, la densidad del combustible
varía entre \num{775} y \SI{840}{\kilogram\per\meter\cubed} con la temperatura
y el origen del crudo \parencite{chevron2007}, y la diferencia entre el poder
calorífico inferior y el superior es del 8 por ciento. El avión reposta en
litros y consume en kilogramos, de modo que el mismo repostaje da un alcance
distinto en un día caluroso que en uno frío. Es un efecto pequeño y es el tipo
de detalle que separa una estimación de una cifra operativa.

\subsection{Razón 7: la distancia tampoco es un número}

Boston y Los Ángeles están a \SI{4193}{\kilo\meter} de círculo máximo, y el
enunciado usa \SI{5000}{\kilo\meter}, un 19 por ciento más. La ruta real es más
larga que el círculo máximo por el encaminamiento del tráfico aéreo, y sobre
todo el alcance efectivo depende del viento: la corriente en chorro sobre
Norteamérica sopla de oeste a este con \num{50} a
\SI{100}{\kilo\meter\per\hour} habituales, lo que en este vuelo, que va contra
ella, equivale a alargar la distancia entre un 8 y un 15 por ciento. La
ecuación de Bréguet trabaja con velocidad respecto al aire y con distancia
respecto al suelo, y quien no separe las dos cosas se equivoca en la dirección
del vuelo.

\subsection{Razón 8: los errores no se suman, se multiplican}

El resultado depende de \(x = Rg/[\eta_o h (L/D)]\) a través de una
exponencial. Para \(x\) pequeño el error relativo se transmite casi
proporcionalmente, pero a medida que \(x\) crece la transmisión se amplifica:
en un vuelo transpacífico con \(x \approx 0{,}8\), un error del 10 por ciento en
el \(L/D\) produce un error del 15 por ciento en el combustible. Y como los
cuatro factores entran multiplicados, un error del 10 por ciento en cada uno de
los tres que se estiman no da un 10 por ciento de error: da un 30 por ciento en
\(x\) si se alinean en el mismo sentido.

\subsection{Razón 9: los datos de entrada no son libres, están acoplados}

Esta es la razón que más me ha sorprendido al hacer el ejercicio y está
desarrollada en la sección~\ref{sec:coherencia}. El enunciado fija cuatro cosas
por separado (peso, velocidad, altitud y \(L/D\)) y esas cuatro cosas están
ligadas por la ecuación de sustentación y por la polar del ala. Al comprobar el
cierre resulta que a Mach 0,948 y \SI{6000}{\meter} un avión de
\SI{10}{\tonne} tiene un \(C_L\) de \num{0.123} y por tanto un \(L/D\) máximo
de \num{5.6}, no de 15. Con el \(L/D\) coherente, el combustible necesario es
de \SI{4958}{\kilogram}, el 50 por ciento del peso al despegue, y el avión no
existe. \textbf{En un problema real no se puede elegir cada dato de entrada por
separado}, y comprobar que el conjunto cierra es tan trabajoso como resolver la
ecuación.

\subsection{Razón 10: lo que se quiere saber no es el queroseno}

Si la pregunta última es la huella climática, el queroseno es un buen
subrogado del \ce{CO2} (la relación es estequiométrica y fija, \num{3.16}) y un
subrogado mediocre del efecto climático total, porque dos tercios de ese efecto
no son \ce{CO2} sino estelas de condensación, óxidos de nitrógeno y vapor de
agua estratosférico \parencite{lee2021}, y esos dos tercios dependen de por
dónde y a qué hora se vuela, no de cuánto se quema. Dos vuelos con el mismo
consumo pueden tener efectos climáticos que difieran en un factor dos según
atraviesen o no una capa de aire sobresaturada de hielo. La ecuación de Bréguet
contesta con precisión a una pregunta que no es exactamente la que interesa.

\begin{table}[htbp]
\small
\caption{Las diez razones, ordenadas por la magnitud del efecto que tienen en
\emph{este} problema. Cálculo propio.}
\label{cuadro:porque}
\begin{tabularx}{\textwidth}{@{}p{0.5cm}Xrp{3.5cm}@{}}
\toprule
\textbf{\#} & \textbf{Razón} & \textbf{Efecto medido aquí} &
\textbf{Naturaleza} \\
\midrule
9 & Los datos de entrada están acoplados por la física y el conjunto no cierra
  & \(+120\) \% & incoherencia del enunciado \\
1 & El combustible pesa: el problema es implícito y hay bola de nieve
  & \(+29\) \% por lectura; \num{0.29} kg por kg de carga
  & estructural, insalvable \\
4 & Reservas normativas de combustible
  & \(+23\) \% & regulatoria \\
5 & \(\eta_o\) y \(L/D\) son los que dominan y son los más difíciles de estimar
  & \(\pm 40\) \% y \(\pm 36\) \% & de modelado \\
7 & La distancia no es la distancia: encaminamiento y viento
  & \(-15\) a \(+9\) \% & de datos \\
2 & La ecuación supone constante lo que no lo es
  & del orden del 5 al 10 \% & de modelado \\
8 & Los errores se multiplican y la exponencial los amplifica
  & \(\times 1{,}5\) a \(x\approx 0{,}8\) & propagación \\
3 & Ascenso y descenso quedan fuera de la ecuación
  & \(+3{,}6\) \% mínimo, el doble en la práctica
  & de alcance del modelo \\
6 & La densidad energética es un intervalo, no un número
  & \(\pm 1\) \% (\(-7\) \% si se confunde con el superior)
  & de datos \\
10 & Lo que se quiere saber es el efecto climático, no el queroseno
  & \(\times 3\) sobre el \ce{CO2} & de definición del problema \\
\bottomrule
\end{tabularx}
\end{table}

\subsection{La respuesta corta}

Si tuviera que contestar la pregunta en dos frases: \textbf{es complicado
porque el combustible es a la vez la respuesta y uno de los datos}, lo que
convierte una división en una ecuación implícita con efecto de realimentación;
y porque \textbf{de los cuatro factores de la ecuación, tres son resultados de
otros tres problemas de ingeniería} (el motor, el ala y la estructura) que sólo
se conocen cuando el avión ya está diseñado. Lo que el ejercicio esconde, y es
lo que a mí me parece la lección de fondo del módulo, es que la ecuación de
Bréguet no es una herramienta de cálculo: es una \textbf{herramienta de
descomposición}. Su valor no está en el número que da, sino en que separa
limpiamente el problema en cuatro figuras de mérito independientes,
\(\eta_o\), \(h\), \(L/D\) y la fracción estructural, y dice exactamente cuánto
vale mejorar cada una \parencite{deweck2022}. Por eso se sigue usando en las
hojas de ruta tecnológicas de la aviación un siglo después de
\textcite{coffin1920}, cuando hace décadas que existen simuladores que dan la
cifra con dos órdenes de magnitud más de precisión.

\section{Respuesta final}

\begin{table}[htbp]
\small
\caption{Respuesta a las dos preguntas del enunciado.}
\label{cuadro:final}
\begin{tabularx}{\textwidth}{@{}p{4.6cm}X@{}}
\toprule
\textbf{Pregunta} & \textbf{Respuesta} \\
\midrule
\textbf{¿Cuánto queroseno hace falta?}
& \textbf{\SI{2248}{\kilogram}}, es decir 2,25 toneladas o unos
  \SI{2800}{\litro}, para el crucero de \SI{5000}{\kilo\meter}. Son el
  \textbf{22,5 por ciento} del peso al despegue, \SI{96.2}{\giga\joule} de
  energía química y \SI{7.1}{\tonne} de \ce{CO2}. Si las \SI{10}{\tonne} se
  interpretan como peso sin combustible, la cifra es \SI{2900}{\kilogram}; con
  la forma lineal de la ecuación, \SI{2547}{\kilogram}; con las reservas
  normativas puestas, \SI{2847}{\kilogram}. La duración del crucero es de
  \SI{4.63}{\hour}. \\
\addlinespace
\textbf{¿Por qué es un cálculo complicado?}
& Porque el combustible pesa y hay que transportarlo, lo que hace el problema
  implícito y logarítmico en lugar de proporcional; porque tres de los cuatro
  factores de la ecuación (\(\eta_o\), \(L/D\) y la fracción estructural) son
  resultados de otros problemas de ingeniería y no datos; porque la ecuación
  supone constante todo lo que en un vuelo real varía; porque la normativa
  añade un 23 por ciento que no está en ninguna ecuación; y porque, como
  muestra la sección~\ref{sec:coherencia}, los datos de entrada están acoplados
  por la física y en este enunciado no cierran: a Mach 0,948 y
  \SI{6000}{\meter} el \(L/D\) alcanzable es \num{5.6} y no 15, con lo que el
  combustible real sería \SI{4958}{\kilogram} y el avión no podría hacer el
  viaje. \\
\bottomrule
\end{tabularx}
\end{table}

\section{Limitaciones}

\begin{enumerate}
\item \textbf{El resultado principal es el que pide el enunciado, no el que
  daría un avión real.} Lo digo en este orden a propósito: la cifra de
  \SI{2248}{\kilogram} es correcta \emph{dados} los datos del enunciado, y los
  datos del enunciado no describen un avión construible
  (sección~\ref{sec:coherencia}). Las dos afirmaciones son compatibles.
\item \textbf{El \(L/D\) alcanzable de \num{5.6} depende de dos supuestos
  míos}: una superficie alar de \SI{29}{\meter\squared}, tomada de un reactor
  de negocios real de masa comparable \parencite{learjet75}, y un coeficiente
  de arrastre parásito de \num{0.022}, que es un valor típico de la literatura
  de diseño \parencite{anderson1999,raymer2018}. Con un ala mayor el \(C_L\)
  baja y el \(L/D\) empeora; con un ala mucho menor la carga alar se vuelve
  impracticable. La conclusión cualitativa es robusta; el número exacto no.
\item \textbf{No he modelado la divergencia de arrastre transónica.} A Mach
  0,948 el arrastre de onda añadiría un incremento de \(C_D\) que no he
  contabilizado, de modo que el \(L/D\) de \num{5.6} es todavía optimista.
\item \textbf{Las reservas de la Figura~\ref{fig:cascada} son un cálculo
  propio con supuestos declarados}: aeropuerto alternativo a
  \num{200} millas náuticas, consumo en espera de
  \SI{400}{\kilogram\per\hour} y rodaje de \SI{40}{\kilogram}. Un despacho real
  usaría las tablas del manual de vuelo del avión y las condiciones
  meteorológicas del día.
\item \textbf{El incremento por ascenso es un mínimo termodinámico}, no una
  estimación operativa. Sólo contabiliza la energía potencial y cinética que hay
  que darle al avión, e ignora que durante el ascenso el \(L/D\) y el
  \(\eta_o\) son peores que en crucero. El valor real está entre dos y tres
  veces por encima.
\item \textbf{La distancia de \SI{4193}{\kilo\meter} es de círculo máximo entre
  las coordenadas de los dos aeropuertos}, cálculo propio, y no una ruta
  publicada. La ruta real es más larga.
\item \textbf{La comparación con el reactor de negocios real usa cifras de
  catálogo} \parencite{learjet75,aopalearjet}, que corresponden a un alcance con
  cuatro pasajeros y reservas incluidas, y por tanto no son estrictamente
  homogéneas con mi cálculo de crucero puro. El factor 1,6 de la comprobación 4
  es un orden de magnitud, no una medida.
\item \textbf{El multiplicador de tres para el efecto climático no \ce{CO2}} es
  el valor central de \textcite{lee2021} para el conjunto de la flota mundial y
  tiene una incertidumbre grande. Aplicado a un vuelo concreto es indicativo, no
  cuantitativo, porque el efecto de las estelas depende de la ruta, la altitud
  y la meteorología de ese vuelo en particular.
\item \textbf{Los factores de emisión no incluyen el ciclo de vida del
  combustible.} Los \SI{3.16}{\kilogram} de \ce{CO2} por kilogramo de queroseno
  son sólo la combustión; la extracción, el refinado y el transporte añaden en
  torno a un 20 por ciento más \parencite{defra2025}.
\item \textbf{La densidad energética que uso, \SI{42.8}{\mega\joule\per\kilogram},
  no es exactamente la del módulo}, que da \SI{42.5}{\mega\joule\per\kilogram}
  siguiendo la tabla de \textcite{tennekes2009}. La diferencia es del 0,6 por
  ciento en el resultado y está recogida en el
  Cuadro~\ref{cuadro:sensibilidad}.
\end{enumerate}

\printbibliography

\end{document}
